% Introduccion
\section{Introduction}
\label{sec:intro}

Before a cargo aircraft takes off, the individual shipments are consolidated onto air cargo
pallets and into containers, the so-called unit load devices (ULDs), which are then loaded into
the aircraft. The shape of a ULD depends on its position in the aircraft: lower-deck containers
have bevelled bottom corners and main-deck pallets and containers have bevelled top corners, so
that they fit the fuselage. Building the ULDs is a daily operational task in air freight
terminals, where it has to be solved for many flights in a short time, and every ULD saved is
space and weight that can be sold or saved in fuel. At the same time, a ULD plan is only usable if it respects the practical
constraints of air transport: the weight capacity of each ULD, the allowed orientations of the
boxes, fragile boxes on which nothing can be placed, the stability of the load, and the weight
distribution, since the center of gravity (CoG) of each ULD must lie in a given area around its
geometric center \citep{Paquay2016}. Building the ULDs is one of the subproblems of the air
cargo load planning problem \citep{Brandt2019}.

Assuming that the cargo consists of rectangular boxes, the problem is a three-dimensional
Multiple Bin Size Bin Packing Problem (MBSBPP) in the typology of \citet{Wascher2007}, with
non-cuboid bins and additional constraints. \citet{Paquay2016} proposed a mixed integer
programming formulation of the problem, which can only solve small instances, and a set of
instances derived from a real set of 562 boxes and the six ULD types of a Boeing 777.
\citet{Paquay2018a} and \citet{Paquay2018b} proposed MIP-based and two-phase constructive
heuristics for these instances. Recently, \citet{Hessler2025} extended the problem with
further practical features (a floor edge for the cargo net locks, on which no box can be loaded,
a substructure to raise the floor, and padding material to support the boxes) and proposed a
randomized greedy search (RGS) based on a fast extreme-point insertion heuristic. The RGS uses
30\% fewer ULDs than the heuristic of \citet{Paquay2018b} in a fraction of a second per
instance, and it is, to the best of our knowledge, the state of the art for these instances.

In this paper we show that a different algorithmic paradigm, borrowed from the single container
loading problem, yields considerably better ULD plans. The beam search guided by a greedy
(BSG) algorithm \citep{Araya2014,Araya2017} loads a container with \emph{blocks} of identical
or similar boxes placed in \emph{maximal free spaces}, and selects the blocks with a beam search
in which each node is evaluated by completing its partial solution with a greedy algorithm. It
is among the best algorithms for the classical container loading problem
\citep{Araya2017,Araya2020}, but the air cargo problem differs from it in several respects:
the ULDs are not cuboids, the support of the boxes and the fragile boxes restrict how boxes can
be stacked, the CoG of every ULD is constrained, several ULDs of different types must be loaded,
and all the boxes must be loaded with as few ULDs as possible. Our contributions are the
following:
\begin{itemize}
\item We adapt BSG to load a ULD with bevelled corners, a floor edge and an optional
	substructure. The blocks are generated so that every box inside a block is supported, and
	each block is placed in a free space only if the boxes of its base are supported according
	to the corner or area criterion of the problem and no box touches a fragile box. A block that
	crosses a bevelled side of the ULD is shifted along the length of the ULD to the nearest
	feasible position.
\item We handle the CoG constraint in four ways: the evaluation of the blocks is penalized by
	the CoG deviation that the placement would cause, the value of a solution is penalized by the
	excess of the CoG deviation over the allowed one, the final load of each ULD is translated as
	a whole to center its CoG, and, when the CoG is still out of range, the ULD is loaded again
	with all the free spaces anchored at the same corner, which groups the load so that it can be
	centered.
\item We embed the single-ULD algorithm into a sequential multi-ULD framework with a
	post-optimization phase that tries to eliminate ULDs by redistributing their boxes into the
	other ULDs and, when the number of ULDs is limited, a ruin-and-recreate phase that inserts the
	boxes left unloaded into the loaded ULDs.
\item We tune the parameters once with irace on a separate training set and evaluate the
	algorithm on the 300 instances of \citet{Paquay2018b} in the three settings considered by
	\citet{Hessler2025}. With the same parameters in all the settings, our algorithm uses
	8--16\% fewer ULDs than the RGS in the settings with an unlimited number of ULDs, where it is
	better in up to 111 of the 300 instances and worse in at most one; with one ULD of each type,
	it leaves fewer boxes unloaded and uses 3--7\% fewer ULDs. Almost none of its ULDs violates the CoG constraint,
	against 7--11\% of the ULDs of the RGS. All the solutions are checked by an independent
	validator, and the code, the solutions and the scripts are publicly available.
\end{itemize}

The remainder of the paper is organized as follows. Section~\ref{sec:problem} defines the
problem and Section~\ref{sec:literature} reviews the related literature.
Section~\ref{sec:algorithm} describes the algorithm, Section~\ref{sec:experiments} reports the
computational experiments, and Section~\ref{sec:conclusions} concludes the paper.
